%% Appendix E ----------------------------------------------------------------
\chapter{Troubleshooting}
\label{app:troubleshooting}
\index{troubleshooting}

\newcommand{\problem}[1]{\par\medskip\noindent{\sffamily\bfseries\color{accentdark}#1}\par\nobreak\smallskip}

\section{Main window}

\problem{Buttons and menu items are grayed out.}
Most analysis controls, and the \gui{Graph Properties} panel, are enabled only after a file is
loaded. \gui{Isolate Burst} additionally requires background subtraction.

\problem{Rectangle zoom does nothing.}
Click \gui{Lock zooming and panning} first, then \gui{Rectangular Zooming}, and drag a
rectangle. Each activation allows one rectangle.

\problem{Drift-rate picking does not stop.}
Right-click or double-click to finish the series of points.

\problem{Selected files cannot be combined.}
The message names the reason --- for example non-consecutive timestamps, mixed stations or an
incomplete time~\(\times\)~focus-code grid. Check the rules in Table~\ref{tab:combine-rules}.

\problem{The UT time axis is unavailable.}
UT display, UT display ranges and UT alignment in the Multi-Station Comparison need the
\texttt{TIME-OBS} keyword in the FITS header. Use seconds instead.

\problem{The thresholds were reset after background subtraction.}
Limits chosen on raw data would saturate the background-subtracted scale and are therefore
reset. Set them again, or apply a preset.

\problem{The application closed unexpectedly.}
Restart it and accept the offer to recover the autosave snapshot, or use \menu{File > Recover
Last Session...}.

\problem{A save fails on Windows.}
The proposed folder may not be writable (for example inside \file{C:\\Program Files}); choose
another folder when asked.

\section{Downloads and archives}

\problem{The station list is empty or a station is missing.}
Station lists contain only stations with data on the selected UTC days. Check the date and the
network connection.

\problem{A listed burst has no FITS files.}
Catalog detections do not guarantee that matching files are archived. Increase
\gui{Extra time}, or search another station.

\problem{Recent SOHO/LASCO or SOHO/EIT data cannot be found.}
The calibrated archives lag real time by many months. Use \gui{Find Latest} to locate the newest
archived data, or \gui{Live Preview (Helioviewer)} for recent quicklook images.

\problem{JSOC downloads, binned frames or cutouts are refused.}
Enter a JSOC notify e-mail registered at
\url{https://jsoc.stanford.edu/ajax/register_email.html} in the \gui{Data Source} card. Binned
and cutout frame sizes require the JSOC source.

\problem{The disk is filling up.}
Empty the caches with the \gui{Clear Cache} buttons of the e-CALLISTO Downloader and the SunPy
Explorer, or with \menu{Data > Reset All (clear cache \& defaults)} in Solar Image Analysis.

\section{Solar Image Analysis and GCS fitting}

\problem{The PFSS card is disabled.}
When running from source, install the \texttt{sunkit-magex} package; the card shows the command.

\problem{Clicked PFSS seeds or the crop rectangle do not respond.}
These interactions need the \gui{PyQtGraph} renderer (\gui{Display \& Crop} card).

\problem{A GCS panel reports that it has no data.}
The STEREO-B archive ends on 27~September~2014. Switch that panel to another channel.

\problem{GCS refinement does not change the direction.}
The direction is constrained only by at least four points in two views separated by
20°--160°. Add points in a second, well-separated view.

\problem{MP4 export is not offered.}
MP4 export needs FFmpeg, which is bundled with the installers. If it is unavailable, save a GIF
instead.

\section{Installation}

\problem{macOS reports that the application cannot be verified.}
Open \menu{System Settings > Privacy \& Security} and click \gui{Open Anyway}, or Control-click
the application and choose \gui{Open}.

\problem{The application does not start on Linux.}
Install the package with \texttt{apt}, which also installs the required Qt/OpenGL libraries
(\texttt{libgl1}, \texttt{libegl1}, \texttt{libxkbcommon-x11-0}, \texttt{libxcb-cursor0}).

\problem{Qt reports ``DLL load failed'' when running from source on Windows.}
Rebuild the virtual environment with \file{packaging/windows/repair_windows_venv.ps1}.

\medskip
If a problem persists, report it with \menu{About > Report a Bug...}
(Section~\ref{sec:bug-report}).
